% erzeugt von 20_code/10_gearcore/scripts/thesis_appendix_tables.py, nicht von Hand ändern
\begin{table}[htbp]
    \caption{Werkstoffklassen des vereinfachten Festigkeitsnachweises nach Stommel et al.: Schubfestigkeitsfaktor $f_{\tau}$ und Steuergröße $q$ \cite[Tabelle 6.1, S.~323]{Sto18}, Abminderungsfaktor $A_{\mathrm{W}}$ der Wechselfestigkeit bei $N = 10^{7}$ \cite[Bild 6.26, S.~362]{Sto18} und Neigungsexponent $k$ der synthetischen Wöhlerlinie \cite[Tabelle 6.3, S.~371]{Sto18}. Die Klasse amorph (verstärkt) trägt für $f_{\tau}$ und $q$ die Werte der spröden Spalte, weil das Buch keine eigene Spalte dafür führt.}
    \label{tab:anhang_klassen_festigkeit}
    \centering
    \footnotesize
    \setlength{\tabcolsep}{4pt}
    \begin{tabular}{L{3.6cm}L{1.9cm}llL{2.4cm}L{2.6cm}}
        Klasse & Verhalten & $f_{\tau}$ & $q$ & $A_{\mathrm{W}}$ & $k$ \\
        \hline
        \arrayrulecolor{black!30}
        teilkristallin (unverstärkt) & duktil & $0{,}58$ & $0{,}00$ & $0{,}20$ bis $0{,}30$ & $11{,}63$ bis $13{,}39$ \\
        \hline
        amorph (unverstärkt) & spröde & $0{,}74$ & $0{,}52$ & $0{,}17$ bis $0{,}20$ & $11{,}63$ bis $12{,}02$ \\
        \hline
        teilkristallin (verstärkt) & semi-duktil & $0{,}65$ & $0{,}26$ & $0{,}25$ & $15{,}65$ bis $20{,}44$ \\
        \hline
        amorph (verstärkt) & spröde & $0{,}74$ & $0{,}52$ & $0{,}17$ & $11{,}39$ bis $13{,}73$ \\
        \hline
        \arrayrulecolor{black}
    \end{tabular}
\end{table}
